% Setup and identification. Statements here; proofs in the appendix.
% Writing rules as in 01_introduction.tex. Numbers only via macros.

\section{Setup and identification}
\label{sec:setup}

\subsection{What is fixed, what is observed, and what is unknown}

The benchmark is fixed. It consists of \NumItems{} items, each with a reference answer, and
for each of \NumModels{} models the leaderboard publishes one response per item. No model is
run again in this paper, and no response is regenerated. Everything below is a statement about
those published responses.

A \emph{verifier} is a function that reads a response and the reference answer and returns
either $1$, meaning the response should be credited, or $0$. A verifier has two parts that
can each fail. It extracts a candidate answer from free text, and it decides whether that
candidate is equivalent to the reference. The first part fails when a response states its
answer in a form the extraction rule does not recognise, and a response whose answer cannot
be extracted is scored $0$, which is the same value a wrong answer receives.

We apply $K = \NumVerifiers$ verifiers to every response, so each response carries a vector of
$K$ verdicts and falls into one of $\NumCells$ \emph{cells}. We write a cell as the string of
verdicts in a fixed order, so that cell $111$ contains the responses every verifier credits
and cell $000$ contains the responses no verifier credits. We call those two cells
\emph{concordant} and the remaining cells \emph{discordant}. Table~\ref{tab:cells} gives the
sizes.

The \emph{human reference} $Y^*$ is defined for every response, although it is observed only
for audited ones. A labeler reads the response, without seeing the reference answer or any
verifier's verdict, and transcribes the final answer the response states, exactly as written.
A labeler who finds no final answer records \textsc{none}, and a labeler who finds two
different final answers both presented as final records \textsc{contradictory}. We then set
$Y^* = 1$ when that transcription matches the reference answer under the benchmark's own
equivalence routine, and $Y^* = 0$ otherwise, counting \textsc{none} and \textsc{contradictory}
as $0$. Section~\ref{sec:sensitivity} reports what changes when they are dropped instead.

Two quantities follow. The first is the \emph{cell share} $s_{mc}$, the fraction of model
$m$'s responses lying in cell $c$. Shares are computed from all \NumRows{} responses and are
known exactly, with no sampling error, because every response is read by every verifier. The
second is $p_{mc}$, the fraction of model $m$'s responses in cell $c$ that the human reference
credits. These are the unknowns. Each model's accuracy under the human reference is
\begin{equation}
A_m \;=\; \sum_{c} s_{mc}\, p_{mc},
\label{eq:accuracy}
\end{equation}
and the quantity a ranking claim is about is the pairwise difference
$\Delta_{jk} = A_j - A_k$.

\subsection{Why the audit alone does not identify the ranking}

Equation~\eqref{eq:accuracy} is linear in the unknown $p_{mc}$, and the coefficients are
known. If every $p_{mc}$ were known, every ordering would follow. An audit estimates them
from a sample, and a sample of a few hundred labels spread over \NumModels{} models and
$\NumCells$ cells leaves most $(m,c)$ pairs with no labels at all.

For a cell with no labels, the only honest statement about $p_{mc}$ is that it lies in
$[0,1]$. Substituting that into~\eqref{eq:accuracy} for both models gives a range for
$\Delta_{jk}$ whose width is the combined share of the unaudited cells. Whenever that width
exceeds the observed gap between two models, the ordering is not identified, and no amount of
statistical machinery applied to the audited cells changes it. This is why the analysis is
organised around cells rather than around models: the discordant cells are small, so bounding
them crudely is affordable, while the concordant cells are large, so they must be bounded
tightly or nothing is identified.

\subsection{The pooling assumption}

Tight bounds on the concordant cells require pooling labels across models, and pooling is an
assumption rather than a measurement. We state it as a single parameter.

\begin{assumption}[Bounded heterogeneity, $H(\varepsilon)$]
\label{ass:h}
For each concordant cell $c$, the values of $p_{gc}$ across the audit sub-populations $g$ that
make up that cell lie in an interval of width at most $\varepsilon$.
\end{assumption}

A sub-population is a model, or a model crossed with a design stratum where the audit design
was stratified more finely. Assumption~\ref{ass:h} restricts how far models may differ from
each other in the cells where the verifiers agree. It says nothing about the level. The
rate at which a human reader overrules a set of agreeing verifiers may be high or low, and the
audit estimates it; $\varepsilon$ bounds only the spread.

Exact homogeneity, $\varepsilon = 0$, is stronger than it looks. Each $p_{gc}$ is a fraction
over finitely many responses, so if the underlying rate is strictly between zero and one, two
sub-populations of finite size will differ by chance alone. $H(0)$ can hold exactly only when
the rate is zero or one. We therefore report certified orderings as a function of
$\varepsilon$, mark the tolerance $\varepsilon = \EpsPrereg$ that we registered before seeing
any result, and treat $\varepsilon = 0$ as the upper end of a curve rather than as an estimate.

\subsection{The identified range}

Let $[L_c, U_c]$ be a confidence interval for the pooled correctness in concordant cell $c$,
constructed in Section~\ref{sec:inference} so that it is valid without assuming the models
share a common rate. For a discordant cell, let $[l_{mc}, u_{mc}]$ be a per-model interval,
which is $[0,1]$ when that cell has no labels for that model.

\begin{proposition}[Identified range]
\label{prop:range}
Fix a pair $(j,k)$ and a concordant cell $c$. Under Assumption~\ref{ass:h} and on the event
$p_{c} \in [L_c, U_c]$, the pair $(p_{jc}, p_{kc})$ lies in
\[
R_c(\varepsilon) \;=\; \bigl\{ (x,y) \in [\ell_c, h_c]^2 \;:\; |x - y| \le \varepsilon \bigr\},
\qquad \ell_c = \max(0, L_c - \varepsilon), \quad h_c = \min(1, U_c + \varepsilon).
\]
Consequently $\Delta_{jk}$ lies in an interval whose endpoints are obtained by minimising and
maximising~\eqref{eq:accuracy} separately in each cell, and for a concordant cell with shares
$a = s_{jc}$ and $b = s_{kc}$ the minimum of $a x - b y$ over $R_c(\varepsilon)$ equals
$a\ell_c - b\min(h_c, \ell_c + \varepsilon)$ when $a \ge b$, and
$a\max(\ell_c, h_c - \varepsilon) - b h_c$ otherwise.
\end{proposition}

At $\varepsilon = 0$ the region collapses to the diagonal, $p_{jc} = p_{kc}$, which is the
pooled analysis; the range is then sharp. For $\varepsilon > 0$ the region is an outer bound,
because the weights with which $j$ and $k$ enter the pooled estimate constrain the attainable
configurations slightly more than $R_c(\varepsilon)$ does. Using the outer bound is
conservative in the direction that matters: it can only widen the reported range, never narrow
it. Appendix~\ref{app:proofs} proves the proposition and the closed form, and
Appendix~\ref{app:tests} checks both against brute-force enumeration.

\subsection{Certification and breakdown}

An ordering $j \succ k$ is \emph{certified} at $\varepsilon$ when the identified range for
$\Delta_{jk}$, widened by the sampling term of Section~\ref{sec:inference}, lies entirely
above zero. Certification is monotone in $\varepsilon$, since $R_c(\varepsilon)$ grows with
$\varepsilon$, which makes the following well defined.

\begin{definition}[Breakdown value]
\label{def:breakdown}
$\varepsilon^*_{jk}$ is the largest $\varepsilon$ at which the ordering $j \succ k$ remains
certified.
\end{definition}

The breakdown value answers the question a skeptical reader should ask, which is not whether
an ordering is certified but how much cross-model heterogeneity it would take to overturn it.
Reporting $\varepsilon^*$ per ordering makes the dependence on Assumption~\ref{ass:h} a
measured quantity rather than a caveat. Because the range in Proposition~\ref{prop:range} is
an outer bound for $\varepsilon > 0$, each $\varepsilon^*_{jk}$ is a lower bound on the true
breakdown value, so the reported robustness is understated rather than overstated.

\subsection{Why more than one verifier is needed}

Stratifying on a single verifier does not help, and the reason is worth stating because it
explains the design.

\begin{proposition}[A common error mechanism cannot reorder]
\label{prop:oneverifier}
Suppose one verifier is used, and suppose its predictive values are common to all models, so
that $\Pr(Y^* = 1 \mid v = 1) = \pi_1$ and $\Pr(Y^* = 1 \mid v = 0) = \pi_0$ regardless of $m$.
Then $\Delta_{jk} = (\pi_1 - \pi_0)(a_j - a_k)$, where $a_m$ is the rate at which the verifier
credits model $m$.
\end{proposition}

Every pairwise difference is scaled by one common factor. A common error mechanism therefore
cannot produce a reordering of one pair while leaving another intact; its only global effect
is to reverse every ordering at once, when $\pi_1 < \pi_0$. Pair-specific reversals, which we
observe, require the mechanism to differ across models. Pooling under a single verifier
assumes away exactly the phenomenon under study, whereas pooling restricted to the cells where
independent verifiers agree leaves the disagreement cells free to be model-specific.
